%% kapa: exam=1 round=32 year=2021 period=2 form=A subject=appraisal_law subject_ko=감정평가관계법규 seq=38 number=38
%% source: Q-Net 감정평가사 기출문제 (cst00309) artlSeq=5208397 file=extracted/1008/5208397/A형/2021년도 제32회 감정평가사 1차 2교시 A형.pdf
%% conversion: kordoc 4.24.51 -> tools/gen1.py (auto); review=proofread-sonnet; flags=none
\begin{question}{38}
도시 및 주거환경정비법령상 조합설립추진위원회(이하 ‘추진위원회’)에 관한 설명으로 옳지 않은 것은?
\begin{choices}
  \item 국토교통부장관은 추진위원회의 공정한 운영을 위하여 추진위원회의 운영규정을 정하여 고시하여야 한다.
  \item 추진위원회는 운영규정에 따라 운영하여야 하며, 토지등소유자는 운영에 필요한 경비를 운영규정에 따라 납부하여야 한다.
  \item 추진위원회는 사용경비를 기재한 회계장부 및 관계 서류를 조합설립인가일부터 30일 이내에 조합에 인계하여야 한다.
  \item 추진위원회는 조합설립에 필요한 동의를 받기 전에 추정분담금 등 대통령령으로 정하는 정보를 토지등소유자에게 제공하여야 한다.
  \item 조합이 시행하는 재건축사업에서 추진위원회가 추진위원회 승인일부터 1년이 되는 날까지 조합설립인가를 신청하지 아니하는 경우에는 정비구역의 지정권자는 정비구역등을 해제하여야 한다.
\end{choices}
\end{question}
\answer{5}
